\subsection*{Auxiliary exact common-$C_3$ comparison after F.10}

This note records an exact model for the residual group in F.7.
It does not assert that the fixed intersection displayed below is trivial.
All permutations use the five-letter identification of $A_5$ already
fixed in F.7--F.10.  Put
\[
 A=\operatorname{Stab}_{A_5}(5),\qquad
 r=(1\ 2\ 3),\qquad K=\langle r\rangle,
 \qquad\beta=(1\ 2)(4\ 5).
\]
Then $A\cong A_4$, $\beta K\beta^{-1}=K$ by inversion,
and $\langle A,\beta\rangle=A_5$.

\paragraph{The actual quotient homomorphism.}
Let $P=P\Gamma_{0,12}$ at the icosahedral configuration.
The discrete branched-cover sequence for the $K$ action is
\[
 1\longrightarrow K\longrightarrow N(K)
       \longrightarrow Q\longrightarrow1,
\]
where $Q$ is the finite-index liftable subgroup of the
mapping class group of the quotient sphere with two decorated
branch points and four free punctures.  The centralizer
$C_P(K)$ has finite index in $N(K)$: both the centralizing
condition and the pure puncture condition have finite index.
Its intersection with $K$ is trivial, since nonidentity deck
maps permute the twelve punctures.  Its image $J$ in $Q$ is
therefore isomorphic to $C_P(K)$ and has finite index in $Q$.
Every element of $J$ fixes the four free quotient punctures.
It also fixes the two branch points separately: exchanging
them conjugates the order-three local deck generator to its
inverse, and hence cannot centralize it.  Thus
\[
             C_P(K)\xrightarrow{\ \sim\ }J
                         \leq P\Gamma_{0,6},
               \qquad [P\Gamma_{0,6}:J]<\infty.          \tag{A.C3.1}
\]
The prime-power comparison and full induced topology proved
in F.7 identify $C_{\widehat P}(K)$ with
$\widehat{C_P(K)}$.  Completing the finite-index embedding
(A.C3.1) identifies it with the specified open subgroup
$\widehat J\leq\widehat{P\Gamma}_{0,6}$.  This gives an
injective continuous homomorphism on the actual completed
centralizer; no completed intersection identity is used.

By F.7.8 the orbit map identifies $C_P(A)$ with the specified
$F_{13}=\pi_1(X)$, where
\[
 X=\mathbf P^1\setminus R_A,\qquad
 R_A=\{\infty\}\cup\{z:z(z^4-1)(z^8+14z^4+1)=0\}.
\]
Since $K\leq A$, composition with (A.C3.1) is the injection
$j:F_{13}\hookrightarrow P\Gamma_{0,6}$ described next.
The solvable comparison and full topology for $A$ and $K$
show that its continuous extension
\[
                \widehat j:\widehat F_{13}
                       \hookrightarrow\widehat{P\Gamma}_{0,6}
                                                               \tag{A.C3.2}
\]
is injective.  This is a different, explicitly specified map
from the thirty-six forgetful maps whose common kernel was
computed in F.9.

\paragraph{The six labelled quotient points.}
In the F.9 tetrahedral coordinate,
\[
 r(z)=\frac{z+i}{z-i},\qquad
 q(z)=z+r(z)+r^2(z)
     =\frac{z^3+3iz+1-i}{(z-1)(z-i)}.
\]
The degree-three quotient $q$ is invariant under $r$.
Its branch points upstairs are
\[
 p_\pm=\frac{(1+i)(1\pm\sqrt3)}2,
\]
and their images are $3p_\pm$.  In the twelve-point
labelling of F.8 the four free $K$-orbits, in order, are
\[
 (0,8,1),\quad(11,3,5),\quad(2,6,9),\quad(7,10,4).
                                                               \tag{A.C3.3}
\]
Their representatives have coordinates $z,1/z,-1/z,-z$.
Consequently the unnormalized labelled quotient map is
\[
 z\longmapsto
   (3p_+,3p_-,q(z),q(1/z),q(-1/z),q(-z)).                \tag{A.C3.4}
\]
Its entries are distinct on $X$: any equality between two
free quotient points would identify their points in the
$A$ orbit, and a collision with a branch point would give a
nontrivial stabilizer in $A$.

Put $U(z)=(z-p_+)/(z-p_-)$.
Direct substitution gives
\[
       \frac{q(z)-3p_+}{q(z)-3p_-}=U(z)^3.
\]
Normalizing the first, second, and third points in (A.C3.4)
to $0,\infty,1$ gives the actual map
\[
 \Sigma:X\longrightarrow M_{0,6},\qquad
 z\longmapsto(s_1(z),s_2(z),s_3(z)),
\]
where
\[
 s_1=\left(\frac{U(1/z)}{U(z)}\right)^3,
 \qquad
 s_2=\left(\frac{U(-1/z)}{U(z)}\right)^3,
 \qquad
 s_3=\left(\frac{U(-z)}{U(z)}\right)^3.                 \tag{A.C3.5}
\]
The homomorphism $j$ in (A.C3.2) is precisely $\Sigma_*$,
by descent of the same labelled orbit path.

\paragraph{The second $A_4$ and the exact based involution.}
The involution $\beta$ exchanges the two branch points and
permutes the four free orbits in (A.C3.3) by $(0\ 1)(2\ 3)$.
On the normalized moduli coordinates its relabelling is
therefore
\[
 B(s_1,s_2,s_3)=(s_1,s_1/s_3,s_1/s_2).                 \tag{A.C3.6}
\]
For example, after swapping the branch points the new
normalizing coordinate is $u\mapsto s_1/u$, since the
first free point has become the old point $s_1$.
This derives (A.C3.6) from the actual labelled quotient,
rather than selecting a conjugate linear action.

Use $z_*$ and the path
$\gamma:z_*\to5+z_*\to b=5+4i$ from (F.10.12).
The point $\Sigma(z_*)$ is fixed by $B$.  Let
$\tau=\Sigma\circ\gamma$, with the indicated direction.
For a chronological loop $\ell$ based at $\Sigma(b)$, the
transported involution is exactly
\[
 \beta_b(\ell)=
 \tau^{-1}*B(\tau)*B(\ell)*B(\tau)^{-1}*\tau.          \tag{A.C3.7}
\]
The path endpoints in this expression match in order;
$B^2=1$ and $B\Sigma(z_*)=\Sigma(z_*)$ show directly that
$\beta_b^2=1$.  It induces the conjugation action of the
same transported finite mapping class $\beta$ used in F.10
on the injective quotient image (A.C3.1).

Since $A$ and $\beta$ generate $A_5$, the exact residual
comparison is
\[
 \widehat j\bigl(C_{\widehat P}(A_5)\bigr)
   =\widehat j(\widehat F_{13})\cap
       \operatorname{Fix}\!\left(
          \beta_b:\widehat{P\Gamma}_{0,6}
                    \longrightarrow\widehat{P\Gamma}_{0,6}
                          \right).                    \tag{A.C3.8}
\]
Indeed centralizing $A_5$ is exactly centralizing $A$ and
$\beta$.  For elements centralizing $A$, equation (A.C3.2)
is injective; also $\beta$ normalizes $K$, so both the
element and its $\beta$ conjugate lie in the domain of
that injection.  Equality of their quotient images is
therefore equivalent to their equality upstairs.

\paragraph{Executed exact checks and their scope.}
The program \path{code/part_a/second_a4_c3_attack.py}
uses exact arithmetic in $\Q(i,\sqrt5,\sqrt{\varphi+2})$
to check the $60$ M\"obius matrices, $180$ generator
products, and $720$ vertex-label images.  It verifies all
four orbits in (A.C3.3), all representative functions,
the quotient identity, and (A.C3.5).  It also checks
the fixed-locus calculation
\[
 \Sigma(z)\in\operatorname{Fix}(B)
 \quad\Longleftrightarrow\quad s_1(z)=s_2(z)s_3(z),
\]
whose numerator, up to a nonzero scalar and after
cancelling the displayed rational denominators, is
\[
 z(z^2+1)(z^8+16z^6+46z^4+16z^2+1).                    \tag{A.C3.9}
\]
It verifies exactly that $z_*$ satisfies the relevant
quartic factor.  Finite intersection of these geometric
loci does not by itself calculate the intersection of
their completed fundamental groups in (A.C3.8).

The auxiliary numerical braid discovery in
\path{work/part_a/f9_c3_johnson_probe.py} currently gives
no exclusion at the first nonzero outer Johnson degree.
It is explicitly exploratory: the crossings there have
not been certified for the actual based map (A.C3.7),
and a higher-degree comparison must retain that path.
No conclusion about (A.C3.8) uses that calculation.
